\scanpage{49}
% L1
provienne d'un $K$-hom $R\mathrel{\hookrightarrow}R'[t_1,\dotsb t_N]$ ($R,R'$ ext. de type fini de $K$), que l'isomorphie du système des $(\gamma_i(R))$ et du système des $(\gamma_i(R'))$ implique que $R$ et $R'$ sont ``équivalents'' en ce sens que chacun s'envoie dans une extension transcendante pure de l'autre \dots

% L9
\uline{Exemples} 1) Soient $C,C'$ deux courbes \inserted{lisses complètes \unsure{conn.|connex.|connexes}} \inserted{sur $k$ alg. clos,} supposons que les jacobiennes soient telles qu'il existe un hom $J\longrightarrow J'$ à noyau fini. Prouver qu'il existe \deleted{$C'$} une application rationnelle \unsure{dominant|dominante}
\[
C'[t_1,\dotsb t_N]\mathrel{\unsure{-|\to}}C\qquad??
\]
% slanted
\margin{\uline{$K$ alg. clos} \uline{au besoin} !}

% L15
2) Si les $\underline{\gamma}^{i}_{\deleted{\unsure{j|i}}}(R)$ sont nuls pour $i\unsure{>| \geq}0$, prouver que $R$ est unirationnelle ?? [En particulier, si la cohomologie \inserted{de $X$ var. proj. non sing.} est entièrement algébrique ??] Cas d'une surface ?? \inserted{[cl. alg. !]}
% archive sheet number: 48
